% An appendix
%======================================================================
\chapter{Wybrane Fragmenty Kodu Źródłowego}
\label{ch:Appendix-Code}
%======================================================================

Niniejszy dodatek zawiera wybrane, kluczowe fragmenty kodu źródłowego projektu \textit{Echoes of Vanitas} z komentarzem objaśniającym ich rolę w architekturze systemu. Pełny kod dostępny jest w repozytorium projektu.

\section{Maszyna Stanów}

Poniższy kod przedstawia implementację kontrolera maszyny stanów oraz klasę bazową stanu postaci.

\begin{verbatim}
// StateMachine.cs
public partial class StateMachine : Node
{
    private CharacterState _currentState;

    public void SwitchState<T>() where T : CharacterState
    {
        _currentState?.ExitState();
        _currentState = GetOrCreate<T>();
        _currentState.EnterState();
    }

    public override void _Process(double delta)
        => _currentState?.UpdateState(delta);

    public override void _PhysicsProcess(double delta)
        => _currentState?.PhysicsUpdateState(delta);
}

// CharacterState.cs - klasa bazowa stanów
public abstract partial class CharacterState : Node
{
    protected Character characterNode;

    public override void _Ready()
        => characterNode = GetOwner<Character>();

    protected virtual void EnterState()  { }
    protected virtual void ExitState()   { }
    public virtual void UpdateState(double delta)       { }
    public virtual void PhysicsUpdateState(double delta){ }
}
\end{verbatim}

\section{Stan Śmierci Przeciwnika}

Fragment implementujący wieloetapową logikę śmierci wroga z obsługą questów i animacji.

\begin{verbatim}
// EnemyDeathState.cs
public partial class EnemyDeathState : EnemyState
{
    private bool _isDying = false;

    protected override void EnterState()
    {
        if (_isDying) return;
        _isDying = true;

        if (characterNode is Enemy enemy)
        {
            enemy.SpawnDrops();
            if (!string.IsNullOrEmpty(enemy.EnemyId))
                QuestManager.Instance?.RegisterKill(enemy.EnemyId);
        }

        characterNode.SetCollisionLayerValue(3, false);
        characterNode.SetCollisionMaskValue(2, false);

        if (characterNode.AnimationPlayerNode.HasAnimation("Death"))
        {
            characterNode.AnimationPlayerNode
                .AnimationFinished += HandleAnimationFinished;
            characterNode.AnimationPlayerNode.Play("Death");
        }
        else
        {
            RemoveEnemy();
        }
    }

    private void HandleAnimationFinished(StringName animName)
    {
        if (animName == "Death")
        {
            characterNode.AnimationPlayerNode
                .AnimationFinished -= HandleAnimationFinished;
            RemoveEnemy();
        }
    }

    private void RemoveEnemy()
    {
        if (characterNode.PathNode != null)
            characterNode.PathNode.QueueFree();
        else
            characterNode.QueueFree();
    }
}
\end{verbatim}

\section{QuestManager – Rejestrowanie Zabójstw}

Fragment systemu questów odpowiedzialny za śledzenie postępu celu Kill.

\begin{verbatim}
// QuestManager.cs
public void RegisterKill(string enemyId)
{
    var questsToComplete = new List<QuestResource>();

    foreach (var quest in ActiveQuests)
    {
        if (quest.Objectives == null) continue;

        foreach (var obj in quest.Objectives)
        {
            if (obj == null) continue;
            if (obj.Type != QuestObjectiveType.Kill) continue;
            if (obj.TargetId != enemyId) continue;

            if (!_killProgress.ContainsKey(quest.Id))
                _killProgress[quest.Id] = new Dictionary<string, int>();
            if (!_killProgress[quest.Id].ContainsKey(enemyId))
                _killProgress[quest.Id][enemyId] = 0;

            _killProgress[quest.Id][enemyId]++;

            if (_killProgress[quest.Id][enemyId] >= obj.RequiredAmount)
                if (!questsToComplete.Contains(quest))
                    questsToComplete.Add(quest);
        }
    }

    foreach (var quest in questsToComplete)
        CompleteQuest(quest);
}
\end{verbatim}

\section{System Dropów Przeciwnika}

Metoda losowania przedmiotów z tablicy dropów wroga z algorytmem Fisher-Yates.

\begin{verbatim}
// Enemy.cs
public void SpawnDrops()
{
    if (DropOrbScene == null || Drops == null || Drops.Length == 0)
        return;

    var rng = new RandomNumberGenerator();
    rng.Randomize();

    var potions = new Array<ItemResource>();
    var others  = new Array<ItemResource>();

    foreach (var it in Drops)
    {
        if (it == null) continue;
        if (it.IsConsumable) potions.Add(it);
        else others.Add(it);
    }

    var droppedOther   = RollOneFromList(others,  rng);
    var droppedPotion  = RollOneFromList(potions, rng);

    if (droppedOther  != null) SpawnOneOrb(sceneRoot, droppedOther,  rng);
    if (droppedPotion != null) SpawnOneOrb(sceneRoot, droppedPotion, rng);
}

// Tasowanie Fisher-Yates + sprawdzenie DropChance
private ItemResource RollOneFromList(Array<ItemResource> list,
                                     RandomNumberGenerator rng)
{
    for (int i = list.Count - 1; i > 0; i--)
    {
        int j = rng.RandiRange(0, i);
        (list[i], list[j]) = (list[j], list[i]);
    }
    foreach (var it in list)
        if (rng.Randf() <= it.DropChance)
            return it;
    return null;
}
\end{verbatim}

\section{Globalny System Zdarzeń}

\begin{verbatim}
// GameEvents.cs
public static class GameEvents
{
    public static event Action<ItemResource, int> OnItemCollected;
    public static event Action<RewardResource>    OnReward;
    public static event Action<int>               OnNewEnemyCount;
    public static event Action                    OnVictory;
    public static event Action                    OnDefeat;

    public static void RaiseItemCollected(ItemResource item, int count)
        => OnItemCollected?.Invoke(item, count);
    public static void RaiseReward(RewardResource r)
        => OnReward?.Invoke(r);
    public static void RaiseNewEnemyCount(int count)
        => OnNewEnemyCount?.Invoke(count);
    public static void RaiseVictory() => OnVictory?.Invoke();
    public static void RaiseDefeat()  => OnDefeat?.Invoke();
}
\end{verbatim}
